% §774 saves `align_state` and drops it to -1000000 *before* `scan_spec` runs,
% and `scan_spec` is what reads `to`/`spread` and the opening brace:
%
%     push_alignment; align_state:=-1000000;
%     @<Scan the preamble and record it in the |preamble| list@>;
%
% The order is not decoration. `scan_keyword` looks one token ahead and puts it
% back when it does not match, and §325's `back_input` *decrements* `align_state`
% for a `{`. In TeX that dip happens down at -1000000 and no one can see it.
% Saving after the keywords captures the dip instead: the `{` that
% `scan_keyword("spread")` pushed back is still pending, so the saved value is -1
% where it should be 0.
%
% §800 hands that -1 back to the enclosing alignment when the inner one ends,
% and `align_state` is no longer 0 there — so §342's test
%
%     if cur_cmd<=car_ret then if cur_cmd>=tab_mark then
%       if align_state=0 then
%       @<Insert the <v_j> template and goto restart@>
%
% never fires, the outer row's `&` stops being a cell boundary, and amsmath's
% `\math@cr@@@ -> &\omit …` arrives at `main_control`: "! Misplaced \omit."
%
% Every alignment nested inside another alignment failed that way — a matrix
% inside an `align*`, which is most mathematical papers. 2608.17265v1 shipped
% zero pages against pdfTeX's eight.
\input plain
\showboxbreadth=100
\showboxdepth=6
% A \halign whose entry contains another \halign. The inner one is the whole
% point; the outer must still see its own `&` and `\cr` afterwards.
\setbox0=\vbox{\halign{#\hfil&#\hfil\cr
  a&\vbox{\halign{#\cr x\cr y\cr}}\cr
  b&c\cr}}
\showbox0
% The same with an explicit `to`, so `scan_keyword("to")` matches instead of
% failing — the branch that does *not* leave a token pushed back.
\setbox0=\vbox{\halign to 100pt{#\hfil&#\hfil\cr
  a&\vbox{\halign{#\cr x\cr y\cr}}\cr
  b&c\cr}}
\showbox0
% `spread`, the other keyword, on the inner alignment.
\setbox0=\vbox{\halign{#\hfil&#\hfil\cr
  a&\vbox{\halign spread 10pt{#\cr x\cr}}\cr
  b&c\cr}}
\showbox0
% Three deep.
\setbox0=\vbox{\halign{#\cr
  \vbox{\halign{#\cr \vbox{\halign{#\cr z\cr}}\cr}}\cr}}
\showbox0
% A \valign inside an \halign entry, so the two directions are mixed.
\setbox0=\vbox{\halign{#\hfil&#\hfil\cr
  a&\hbox{\valign{#\cr\hbox{x}\cr\hbox{y}\cr}}\cr
  b&c\cr}}
\showbox0
\end
